% 349_SU2L_Chiralitaet_Galois_En_ch.tex
\providecommand{\BM}{\textbf{[B]}}\providecommand{\KM}{\textbf{[K]}}\providecommand{\SM}{\textbf{[S]}}
\chapter*{SU$(2)_L$ Chirality from the Galois Structure:\\
Su/Sd Split and Cl$(6)$ Grade Parity}
\addcontentsline{toc}{chapter}{SU$(2)_L$ Chirality from Galois}
\begin{center}\large\textbf{Johann Pascher}\\\normalsize ORCID 0009-0000-6518-4064\\4~September 2026\end{center}
\section*{Status markers}\begin{center}\begin{tabular}{@{}lp{10.5cm}@{}}\textbf{[B]} & \emph{Established} --- secured by script.\end{tabular}\end{center}
\section*{Abstract}
The SU$(2)_L$ chirality structure of the Standard Model follows algebraically from the Cl$(6)$ grade parity established in Doc.~341. The pseudoscalar $\gamma_5=e_1e_2e_3e_4e_5e_6$ has eigenvalue $(-1)^r$ on grade-$r$ elements. The Su/Sd split (Doc.~344) is therefore exactly the right/left split: Su = even grades = right-handed \BM, Sd = odd grades = left-handed \BM. The open question from Docs.~347/348 is thereby closed \BM.
\section{Theorem A --- Grade parity and $\gamma_5$ eigenvalue \BM}
$\gamma_5\cdot X=(-1)^r X$ for grade-$r$ element $X\in\mathrm{Cl}(6)$. Even grades: $\gamma_5=+1$ (right-handed). Odd grades: $\gamma_5=-1$ (left-handed). \BM
\section{Theorem B --- Su = even grades = right-handed \BM}
Su[0]=$\nu$: grade 0 $\to$ right. Su[1]=$\bar u_d$: grade 2 $\to$ right. Su[2]=$u$: grade 4 $\to$ right. Su[3]=$e_R$: grade 6 $\to$ right. All $\gamma_5=+1$. \BM
\section{Theorem C --- Sd = odd grades = left-handed \BM}
Sd[0]=$\bar\nu$: grade 1 $\to$ left. Sd[1]=$d$: grade 3 $\to$ left. Sd[2]=$\bar d$: grade 5 $\to$ left. Sd[3]=$e_L$: grade 7 $\to$ left. All $\gamma_5=-1$. \BM
\section{Theorem D --- SU$(2)_L$ doublets from Sd sector \BM}
Left-handed doublets: $(\nu_L,e_L)=(\mathrm{Sd}[0],\mathrm{Sd}[3])$; $(u_L,d_L)=(\mathrm{Sd\text{-quarks}})$. Right-handed singlets: $e_R=\mathrm{Su}[3]$; $u_R=\mathrm{Su}[2]$; $\nu_R=\mathrm{Su}[0]$. SU$(2)_L$ acts on the Sd sector (left-handed) only. \BM
\section{Theorem E --- $\mathbb{Z}_3$ fixed points/orbit = chirality \BM}
From Doc.~341: $\mathrm{Vss}[0,2,4]$ ($\sigma$-fixed points) = right-handed; $\mathrm{Vss}[1,3,5]$ ($\sigma$-three-orbit) = left-handed. The SU$(2)_L$ chirality split is a direct consequence of the $\mathbb{Z}_3$ symmetry of T$^4$/Z$_3$. \BM
\section{Summary}
\begin{table}[H]\centering\begin{tabular}{lll}\toprule Theorem & Content & Status\\\midrule
A & $\gamma_5$ eigenvalue $(-1)^r$ on grade $r$ & \BM\\
B & Su = even grades = right-handed & \BM\\
C & Sd = odd grades = left-handed & \BM\\
D & SU$(2)_L$ doublets = Sd sector & \BM\\
E & $\mathbb{Z}_3$ fixed points/orbit = right/left & \BM\\\bottomrule\end{tabular}\end{table}
The SU$(2)_L$ chirality structure is now derived from the Galois/Clifford algebra \BM. Open: exact CKM mixing angles and CP phase $\delta_{CP}$ (Doc.~348) \SM.
\section*{Verification script}\texttt{pruef\_349\_chiralitaet.py} --- Theorems A--E, all assertions passed \BM.
\section*{Note on AI assistance}Prepared with support from Claude (Anthropic). Physical interpretations: Johann Pascher (ORCID 0009-0000-6518-4064).
\begin{thebibliography}{9}
\bibitem{P341} J.~Pascher, Doc.~341, FFGFT corpus (Sept.\ 2026).
\bibitem{P344} J.~Pascher, Doc.~344, FFGFT corpus (Sept.\ 2026).
\bibitem{P347} J.~Pascher, Doc.~347, FFGFT corpus (Sept.\ 2026).
\bibitem{Furey2015} C.~Furey, Ph.D.\ Thesis, Waterloo 2015, arXiv:1611.09182.
\end{thebibliography}